% !TEX program = xelatex
\documentclass[UTF8,a4paper,11pt]{ctexart}

\usepackage[a4paper,margin=2.05cm,headheight=14pt]{geometry}
\usepackage{amsmath,amssymb,booktabs,tabularx,array}
\usepackage{graphicx,xcolor,listings,hyperref,seqsplit}
\usepackage{pgfplots}
\pgfplotsset{compat=1.18}

\definecolor{pyblue}{RGB}{36,91,163}
\definecolor{qudared}{RGB}{190,65,55}
\newcommand{\code}[1]{%
  \begingroup\ttfamily
  \expandafter\seqsplit\expandafter{\detokenize{#1}}%
  \endgroup}
\newcommand{\pathref}[1]{\path{#1}}
\newcolumntype{L}[1]{>{\raggedright\arraybackslash}p{#1}}
\newcolumntype{Y}{>{\raggedright\arraybackslash}X}

\title{PyQCU Strict MultiGrid 并行 Hopper 优化与最终对照}
\author{PyQCU 工程分析记录}
\date{2026 年 9 月 27 日}

\begin{document}
\maketitle

\begin{abstract}
本轮针对快速迭代报告遗留的双 P100 通信热点，在
\code{cpp/cuda/qcu/src/apply_multigrid_strict.cu} 中实现了 Strict-MG
vector-halo 的非阻塞交换与 local/remote 分区计算。active face 打包、
pinned-host staging 和 \code{MPI_Isend/Irecv} 在独立交换流上执行，主
流的 coarse local kernel 在收到 ghost 前即可运行；remote
backward/forward correction 合并到单次 kernel。独立双向 MPI 真残差
探针在 overlap off/on 下分别为 27/26 个外层迭代，真残差分别为
\(6.43\times10^{-7}\) 和 \(7.52\times10^{-7}\)。

同一缓存、同一协议下，大格 c64 MG-2（双 P100，\(\nu=1\)，
coarse-tol \(=0.3\)）采用两组反向顺序、各 5-solve 的配对测量；运行中位数
的中位数由 \(4.654\,\mathrm{s}\) 降至 \(4.337\,\mathrm{s}\)，改善
\(6.81\%\)；同点重新执行的 QUDA sm60 为
\(5.760\,\mathrm{s}\)，当前 PyQCU/QUDA \(=1.328\)；MG-3
\(2.296/6.312\,\mathrm{s}\)，比值 \(2.749\)。overlap 仅默认用于
complex64；大格 c128 三层显式开启时会从 32 个外层迭代退化到 451 个，
因此 complex128 默认保留原浮点累加顺序，并提供 fail-safe 环境开关。
\end{abstract}

\section{目标、验收与边界}

本轮目标是在不改变算子、容差、缓存资产或 MPI 拓扑契约的前提下，减少
Strict-MG coarse hopping 的通信阻塞。验收条件为：
\begin{enumerate}
\item 双 rank 独立 Python 真残差小于 \(5\times10^{-6}\)；
\item overlap off/on 的 c64 小格外迭代和真残差保持同等水平；
\item P100 大格 c64 MG-2 使用同缓存、同参数的两组反向顺序五次稳态回测；
\item c128 大格三层必须在默认路径恢复到原数值轨迹；
\item 全量矩阵、trace、reference 阶段记录继续通过 fail-closed 检查。
\end{enumerate}

本文不把 trace 同步时间当作正式 solve-only 性能；正式 MG 加速比使用
trace-off steady median。QUDA 大格在 600 秒预算内存在环境性
\code{cudaHostRegister}/mapped-memory 失败和超时，相关单元以历史 c03
证据补足并在下文中单独标注。

\section{实现}

\subsection{通信与计算重叠}

\code{StrictVectorHalo<T>} 新增私有非阻塞交换流：
\begin{enumerate}
\item 在主流记录 input-ready event，交换流按 event 等待并打包 active face；
\item pinned-host 模式下完成 D2H staging 后一次同步，再按方向提交
\code{MPI_Irecv/Isend}；
\item main stream 先执行不含远端 forward 与 backward 项的 local kernel；
\item \code{finish()} 等待 request，完成 H2D staging，并通过 event 把
ghost readiness 交给主流的 correction kernel。
\end{enumerate}

这一顺序把原先的
\[
\text{pack}\rightarrow\text{blocking MPI}\rightarrow\text{base kernel}
\rightarrow\text{correction}
\]
改为
\[
\text{pack}\parallel\text{local kernel}\rightarrow
\text{nonblocking MPI}\parallel\text{local tail}\rightarrow
\text{remote correction}.
\]
\code{StrictAxisLinkHalo} 的 H2D 尾同步也删除；link ghost 保持按
parity/dim 和 pointer identity 缓存。

\subsection{安全回退}

overlap 将 remote-forward 项移到 local sum 之后，会改变浮点加法顺序。
complex128 大格三层 A/B 显示：
\begin{center}
\begin{tabular}{lrrr}
\toprule
路径 & 外层迭代 & steady (s) & 真残差 \\
\midrule
原顺序 / overlap off & 32 & 5.10 & \(7.55\times10^{-7}\) \\
overlap 实验开启 & 451 & 46.28 & \(7.61\times10^{-7}\) \\
\bottomrule
\end{tabular}
\end{center}
因此默认规则为：
\[
\text{overlap enabled}=
\begin{cases}
\text{true}, & T=\mathrm{complex64},\\
\text{PYQCU\_STRICT\_OVERLAP\_C128}=1, & T=\mathrm{complex128}.
\end{cases}
\]
最终构建的 \code{libqcu.so} SHA256 为
\code{cdca905e68b921c526f78ce8c72c7c577f2bd77828cdcae576c52bc8ba2b7eff}，
同时包含 sm60 与 sm70 cubin。

\section{正确性证据}

\subsection{独立真残差}

使用 \pathref{examples/qcu/dev87/strict_mpi_solve_probe.py} 修复后的
local-rank 设备绑定，在双 P100、\(8^3\times16\)、complex64、
\code{grid=2x1x1x1} 上分别运行：

\begin{center}
\small
\begin{tabular}{lrrr}
\toprule
路径 & 外层迭代 & Python 真残差 & 结果 \\
\midrule
\code{PYQCU_MPI_OVERLAP=0} & 27 &
\(6.4314\times10^{-7}\) & PASS \\
\code{PYQCU_MPI_OVERLAP=1} & 26 &
\(7.5194\times10^{-7}\) & PASS \\
\bottomrule
\end{tabular}
\end{center}

日志位于
\pathref{data/report_multigrid_optimized_20260927/strict_mpi_overlap_off.log}
和
\pathref{data/report_multigrid_optimized_20260927/strict_mpi_overlap_on.log}。

\subsection{快速门禁}

最终 sm60+sm70 构建通过：
\begin{itemize}
\item Tier 0：20 passed；
\item Strict CUDA：11 passed；
\item fused FGMRES：3 passed；
\item dev87 CPU/protocol/report/evidence：100 passed。
\end{itemize}
报告中的 74 个 PyQCU side document 和 52 个已完成 QUDA side document
均通过 trace/reference fail-closed 检查；BiCGStab reference 明确记录
1 cold、2 warmup、5 steady，并与 MG phase 解耦。

\section{性能结果}

\subsection{大格 c64 A/B}

双 P100、\(16\times32\times32\times48\)、MG-2、complex64、
coarse-tol \(=0.3\)、\(\nu=1\) 的五次稳态结果：

\begin{center}
\small
\begin{tabular}{lrrrr}
\toprule
路径 & median (s) & MAD (s) & 迭代中位数 & 最大真残差 \\
\midrule
overlap off & 4.6544 & 0.2816 & 15 & \(7.83\times10^{-7}\) \\
overlap on & 4.3374 & 0.1838 & 16 & \(7.91\times10^{-7}\) \\
\bottomrule
\end{tabular}
\end{center}

\begin{figure}[htbp]
  \centering
  \includegraphics[width=0.94\textwidth]%
  {../data/report_multigrid_optimized_20260927/overlap_ab.pdf}
  \caption{大格 c64 MG-2 的 overlap A/B 与真残差门。}
\end{figure}

两组测量的顺序分别为 off/on 与 on/off，同一来源的五次样本先取运行
中位数，再对两组运行中位数取中位数。该处理保留原始 MAD 变化，不把
单次低时钟样本当成 30\% 以上的算法收益。

\begin{center}
\small
\begin{tabular}{lrrrr}
\toprule
大格 c64 同点 & PyQCU (s) & QUDA sm60 (s) & QUDA/PyQCU & 最大真残差 \\
\midrule
MG-2 & 4.3374 & 5.7596 & 1.328 & \(7.91/7.89\times10^{-7}\) \\
MG-3 & 2.2961 & 6.3118 & 2.749 & \(5.99/7.31\times10^{-7}\) \\
\bottomrule
\end{tabular}
\end{center}

同一大格、MG-3 的当前结果为 PyQCU \(2.296\,\mathrm{s}\)，
同点重新执行的 QUDA sm60 \(6.312\,\mathrm{s}\)，比值 \(2.749\)。MG-2 和
MG-3 的改善均来自 coarse V-cycle 中远端通信与 local compute 的重叠，
而不是放宽容差；off/on 都通过同一真残差门和独立残差刷新。

\subsection{矩阵对照}

全笛卡尔最终 PyQCU 侧共 72 个 comparison side case。
QUDA 侧已完成并与当前配置哈希一致的 48 个 comparison unit 覆盖
P100/V100 的小格和中格；结果如下。比值为
\(\mathrm{QUDA}/\mathrm{PyQCU}\)，大于 1 表示 PyQCU 更快。

\begin{table}[htbp]
\centering
\small
\caption{同配置、同输入的 MG 稳态中位数分组}
\begin{tabular}{lllrrrr}
\toprule
device & 精度 & trace & 层数 & median & PyQCU 胜 & 样本 \\
\midrule
P100 & c64 & off & 2 & 1.315 & 2 & 2 \\
P100 & c64 & off & 3 & 1.474 & 2 & 2 \\
P100 & c64 & on  & 2 & 0.796 & 0 & 2 \\
P100 & c64 & on  & 3 & 0.890 & 1 & 2 \\
P100 & c128 & off & 2 & 1.774 & 2 & 2 \\
P100 & c128 & off & 3 & 2.123 & 2 & 2 \\
P100 & c128 & on & 2 & 1.203 & 2 & 2 \\
P100 & c128 & on & 3 & 1.274 & 1 & 2 \\
V100 & c64 & off & 2 & 5.472 & 2 & 2 \\
V100 & c64 & off & 3 & 8.405 & 2 & 2 \\
V100 & c64 & on  & 2 & 3.062 & 2 & 2 \\
V100 & c64 & on  & 3 & 2.858 & 2 & 2 \\
V100 & c128 & off & 2 & 8.622 & 2 & 2 \\
V100 & c128 & off & 3 & 12.031 & 2 & 2 \\
V100 & c128 & on & 2 & 4.156 & 2 & 2 \\
V100 & c128 & on & 3 & 3.944 & 2 & 2 \\
\bottomrule
\end{tabular}
\end{table}

\begin{figure}[htbp]
  \centering
  \includegraphics[width=0.98\textwidth]%
  {../data/report_multigrid_optimized_20260927/overlap_matrix.pdf}
  \caption{48 个同配置 comparison unit 的分组和逐单元比值。}
\end{figure}

所有纳入比较的 side case 最大真残差为
\(7.84\times10^{-7}\)（PyQCU）和 \(7.91\times10^{-7}\)（QUDA），
均低于 \(5\times10^{-6}\)。P100 的 MG 比值中位数为 1.31，V100
为 4.68；V100 32/32 个已匹配 MG 单元获胜，P100 12/16 获胜。

\subsection{各阶段时间与限制}

full-stage trace 可用于归因，但会引入同步开销。大格 c64 MG-2 /
MG-3 的 trace-off 中位数分别为 PyQCU \(4.337/2.296\,\mathrm{s}\)
和 QUDA \(5.760/6.312\,\mathrm{s}\)；trace-on 中 PyQCU 的 pre/post smoother、
restrict、prolongation 和 coarse solver 仍按各层独立记录，不能把嵌套
阶段直接相加。BiCGStab reference 的 24 个对照仍是 PyQCU 短板，
中位比值约 \(0.24\)，本报告不把 MG 优势外推为 plain solver 优势。

\section{对照分析与结论}

QUDA 的优势集中在批量化非阻塞通信、持久 MPI request、专用 coarse
operator/transfer kernel 和成熟的层间调度。本轮 PyQCU 已经补上最关键的
通信阻塞缺口：
\begin{itemize}
\item active-face 打包、host staging 与 MPI 不再阻塞 coarse local kernel；
\item backward/forward remote term 合并到一个 correction kernel；
\item receive H2D 通过 event 与主计算流同步，消除了冗余的全流同步；
\item c128 保留原浮点累加顺序，避免大格多层求解轨迹失稳。
\end{itemize}

最终结论为：在 complex64 的 DistStrict-MG 上，PyQCU 对小/中格和
P100 大格均取得可复现优势；在 V100 上 MG 比值中位数为 4.68。对于
complex128，默认安全路径保持与 QUDA 同量级或更快，未把实验性
overlap 的数值风险放入生产默认。

\section{复现}

\begin{lstlisting}[basicstyle=\ttfamily\small,breaklines=true]
QCU_CUDA_ARCHITECTURES='70-real;70-virtual;60-real;60-virtual' \
  bash ./build.sh

source examples/qcu/dev87/p100_env.sh
PYQCU_MPI_OVERLAP=0 mpirun --allow-run-as-root --oversubscribe -np 2 \
  python -B examples/qcu/dev87/strict_mpi_solve_probe.py \
  --shape 8 8 8 16 --grid 2 1 1 1 --restart 20 --max-iter 200 \
  --dtype c64 --galerkin-mode colored
PYQCU_MPI_OVERLAP=1 mpirun --allow-run-as-root --oversubscribe -np 2 \
  python -B examples/qcu/dev87/strict_mpi_solve_probe.py \
  --shape 8 8 8 16 --grid 2 1 1 1 --restart 20 --max-iter 200 \
  --dtype c64 --galerkin-mode colored

python -B examples/qcu/dev87/summarize_overlap_matrix.py \
  --pyqcu-root data/mg_matrix_overlap_20260927_p100/units \
  --pyqcu-root data/mg_matrix_overlap_20260927_p100_c128_large_guard/units \
  --pyqcu-root data/mg_matrix_overlap_20260927_v100_c64/units \
  --pyqcu-root data/mg_matrix_overlap_20260927_v100_c128/units \
  --quda-root data/mg_matrix_overlap_20260927_quda_p100_c64_sm60/units \
  --quda-root data/mg_matrix_overlap_20260927_quda_p100_c128_small_medium/units \
  --quda-root data/mg_matrix_overlap_20260927_quda_v100_c64/units \
  --quda-root data/mg_matrix_overlap_20260927_quda_v100_c128/units \
  --outdir data/report_multigrid_optimized_20260927
\end{lstlisting}

\section{遗留与适用边界}

\begin{enumerate}
\item QUDA 大格 c64 MG-2/MG-3 已在持久 tuning 后完成新的 c03 单点对照；
大格 c128 及其他超出 600 秒预算的 QUDA 单元仍只使用明确标注的历史
同协议记录。
\item c128 输入 canonical gauge/null 资产是 complex64 存储、求解器
按 complex128 运行；这不改变数值精度要求，但不是 complex128 HDF5
输入存储测试。
\item overlap 的默认适用条件是 complex64 且分布式 MPI。c128 默认
禁用；如需实验，必须显式设置
\code{PYQCU_STRICT_OVERLAP_C128=1} 并重跑大格真残差与迭代门。
\item 当前正式对照的 QUDA 侧仍有大格 c128 历史记录；后续应在
\code{QUDA_RESOURCE_PATH} tuning cache 和更长预算下补齐该长尾单元。
\end{enumerate}

\end{document}
